% --- Back cover ---
% Large PERSONIX wordmark, no logo image. Cream paper inherited globally.
% Absolute (tikz overlay) positioning so the footer sticks to the bottom.
\clearpage
\thispagestyle{empty}
\begin{tikzpicture}[remember picture, overlay]
  % --- Wordmark ---
  \node[anchor=north, font=\sffamily\bfseries\fontsize{48}{52}\selectfont]
    at ([yshift=-26mm]current page.north) {PERSONIX};
  \node[anchor=north, font=\rmfamily\itshape\large]
    at ([yshift=-44mm]current page.north) {Kompromissloser Wandel};
  \node[anchor=north, font=\rmfamily\small,
        text width=\paperwidth-50mm, align=center]
    at ([yshift=-54mm]current page.north)
    {Ein unzensierbares und unbestechliches dezentrales Reputationsnetzwerk};
  % --- Body ---
  \node[anchor=north, text width=\paperwidth-50mm, align=justify, font=\small]
    at ([yshift=-72mm]current page.north)
    {%
      Der Staat funktionierte, als die Kommunikation langsam war, Entscheidungen
      lokal fielen und die Autorität in derselben Stadt lebte, die sie regierte.
      Die heutigen Netzwerke kehren alle drei Annahmen um --- und die auf ihnen
      errichteten Institutionen passen nicht mehr zu der Welt, die sie regieren.
      Dieses Buch beschreibt ein Werkzeug, das passt:
      ein dezentrales Reputationsnetzwerk, in dem die Identität dir gehört,
      die Wahrheit öffentlich prüfbar ist und Bestechung reputationeller
      Selbstmord ist.

      \smallskip
      Kein Manifest. Keine Prognose. Ein funktionierender Bauplan dafür, wie sich
      der Nachfolger des Staates errichten lässt --- freiwillig, eine Erklärung
      nach der anderen.
    };
  % --- Footer (anchored to the bottom of the page) ---
  \node[anchor=south, font=\sffamily\footnotesize,
        text width=\paperwidth-50mm, align=center]
    at ([yshift=12mm]current page.south)
    {%
      Pavel Kudrna\quad\textbullet\quad Prag, 2026\par
      \href{https://personix.org}{personix.org}\par
      Lizenziert unter Creative Commons Attribution-ShareAlike 4.0 International
      (CC BY-SA 4.0).
    };
\end{tikzpicture}%
\mbox{}%
\clearpage
